\section*{\centering Глава 1. Аналитическая часть}
\addcontentsline{toc}{section}{ Глава 1. Аналитическая часть}

Не сложно понять желание людей найти легкий способ обогатиться. Некоторые готовы переступить через моральные принципы и не прочь обманывать тех, кого несложно обмануть. Для них обман -- это инструмент, с помощью которого можно получить деньги, ресурсы или выгоду, не прилагая усилий и не сталкиваясь с последствиями. Но любому действию есть противодействие, поэтому различные институты -- от государственных до частных -- пытаются противостоять этому. Из-за этого мошенники в постоянном поиске новых способов обмана. 

С развитием технологий стали появляться новые виды мошенничества. Злоумышленники активно используют современные средства связи, социальные сети и цифровые платформы для реализации своих схем. Среди наиболее распространенных методов можно выделить фишинг, кражу персональных данных\footnote{Кража персональных данных -- это незаконное получение и использование личной информации жертв для финансовой выгоды или манипуляции.}, поддельные онлайн-магазины\footnote{Поддельные онлайн-магазины -- это фальшивые веб-сайты, создаваемые мошенниками для обмана покупателей и получения денег за товары, которые никогда не будут доставлены.} и телефонное мошенничество.

Телефонное мошенничество часто основывается на непосредственном контакте с жертвой, что делает его более личным и убедительным. Мошенники используют доверие и страх, чтобы манипулировать людьми. Эти схемы могут быть более сложными, поскольку включают элементы эмоционального воздействия и срочности.

Для борьбы с такими методами необходимо создать систему, которая будет включать не только повышение осведомленности граждан, но и технологии для проверки и блокировки подозрительных звонков. Важным аспектом является также обучение людей распознавать мошенников и активное взаимодействие с правоохранительными органами для пресечения этих схем.

В данной главе будут рассмотрены подходы к решению задач работы, а именно методы определения телефонного мошенничества и организации реестра телефонных мошенников. Будет проведен их обзор, классификация и сравнение, а также рассмотрены положительные и отрицательные аспекты существующих решений. Определены требования к нашему решению и составлено техническое задание.

\subsection*{1.1. Методы определения телефонного мошенничества}
\addcontentsline{toc}{subsection}{1.1. Методы определения телефонного мошенничества}

- перечислить все методы определения телефонного мошенничества | три страницы
- классифицировать эти методы | одна страница
- определить плюсы и минусы этих методов | одна страница

/* на пять страниц */

\subsection*{1.2. Помощник по обнаружению мошенничества}
\addcontentsline{toc}{subsection}{1.2. Помощник по обнаружению мошенничества}

- рассказать про Spam Detection Assistant

/* на четыре страницы */

\subsection*{1.3. Реестр телефонных мошенником}
\addcontentsline{toc}{subsection}{1.3. Реестр телефонных мошенником}

- найти методы организации реестра телефонных мошенников
- классифицировать эти методы
- определить плюсы и минусы этих методов

/* на четыре страницы */

- представить наше решение: обнаружение мошенничества + реестр мошенников + мобильное приложение
- техническое задание, требования к аппаратным средствам, необходимым для работы системы

/* на одну страницу */

\newpage
